\chapter{Thực nghiệm và Kết quả}
\label{ch:experiment}

Chương này trình bày quá trình đánh giá hệ thống. Trước hết, Mục~\ref{sec:data-stats} thống kê bộ
dữ liệu CronQ-VN. Mục~\ref{sec:setup} mô tả thiết lập thực nghiệm, Mục~\ref{sec:baselines} giới
thiệu các phương pháp so sánh và độ đo. Các mục tiếp theo lần lượt trả lời ba câu hỏi nghiên cứu:
kết quả tổng (Mục~\ref{sec:overall-results}, CHNC1), nghiên cứu cắt bỏ (Mục~\ref{sec:ablation},
CHNC2) và phân tích bộ lọc cross-encoder (Mục~\ref{sec:ce-analysis}, CHNC3). Cuối cùng,
Mục~\ref{sec:error-analysis} và Mục~\ref{sec:discussion} phân tích hiệu quả suy luận, lỗi điển hình
và bàn luận hạn chế.

\medskip
\noindent\textbf{Lưu ý:} tại thời điểm hoàn thiện bản thảo này, pha đánh giá định lượng đang được
tiến hành. Các bảng kết quả ở Mục~\ref{sec:overall-results}--\ref{sec:ce-analysis} đã được dựng sẵn
khung; các ô số liệu đánh dấu ``--'' sẽ được điền sau khi hoàn tất chạy đánh giá.

\section{Bộ dữ liệu và thống kê}
\label{sec:data-stats}

Đồ thị tri thức thời gian nền của CronQ-VN, sau khi trích xuất và lọc từ Wikidata
(Mục~\ref{sec:buildkg}), có quy mô tóm tắt trong Bảng~\ref{tab:kg-stats}.

\begin{table}[ht]
\centering
\caption{Thống kê đồ thị tri thức thời gian nền của CronQ-VN.}
\label{tab:kg-stats}
\begin{tabular}{L{7.0cm} R{3.5cm}}
\toprule
\textbf{Thành phần} & \textbf{Số lượng} \\
\midrule
Sự kiện thời gian (bộ bốn) & $\sim$237.900 \\
Thực thể có nhãn tiếng Việt & $\sim$37.580 \\
Quan hệ & 271 \\
Quan hệ dùng để sinh câu hỏi & 21 \\
\bottomrule
\end{tabular}
\end{table}

Từ đồ thị này, bộ câu hỏi CronQ-VN được sinh ra theo năm loại với phân bố mục tiêu đồng đều (20\%
mỗi loại). Bảng~\ref{tab:q-stats} trình bày phân bố câu hỏi theo loại. Do bốn trong năm loại câu hỏi
có đáp án là thực thể và chỉ một loại (\texttt{simple\_time}) có đáp án là thời gian, tập dữ liệu
nghiêng về đáp án dạng thực thể -- một đặc điểm cần lưu ý khi phân tích kết quả theo kiểu đáp án.

\begin{table}[ht]
\centering
\caption{Phân bố câu hỏi trong CronQ-VN theo loại.}
\label{tab:q-stats}
\begin{tabular}{L{4.5cm} C{3cm} C{3cm}}
\toprule
\textbf{Loại câu hỏi} & \textbf{Kiểu đáp án} & \textbf{Số câu} \\
\midrule
\texttt{simple\_entity} & Thực thể & -- \\
\texttt{simple\_time}   & Thời gian & -- \\
\texttt{before\_after}  & Thực thể & -- \\
\texttt{first\_last}    & Thực thể & -- \\
\texttt{time\_join}     & Thực thể & -- \\
\midrule
\textbf{Tổng} & & \todo{tổng số câu sau khi sinh đủ} \\
\bottomrule
\end{tabular}
\end{table}

\noindent\todo{điền số câu mỗi loại sau khi sinh xong bộ dữ liệu đầy đủ (mục tiêu 30.000 câu).}

\section{Thiết lập thực nghiệm}
\label{sec:setup}

Toàn bộ pha suy luận được vận hành cục bộ. Bảng~\ref{tab:setup} liệt kê các siêu tham số và cấu
hình mô hình, đặt theo đúng tinh thần §4.1 của~\cite{chen2024ari}.

\begin{table}[ht]
\centering
\caption{Cấu hình thực nghiệm.}
\label{tab:setup}
\begin{tabular}{L{7.5cm} L{5.0cm}}
\toprule
\textbf{Thành phần / siêu tham số} & \textbf{Giá trị} \\
\midrule
Mô hình ngôn ngữ (suy luận / pha đánh giá) & Qwen (\texttt{qwen3:14b}, qua Ollama) \\
Mô hình ngôn ngữ (chắt lọc methodology / pha học) & GPT-4o (\texttt{gpt-4o}, theo bản gốc) \\
Mô hình nhúng & \texttt{nomic-embed-text} \\
Nhiệt độ giải mã $T$ & 0 (tất định) \\
Số mẫu lịch sử để học & 200 \\
Số cụm K-means $k$ & 10 \\
Số bước suy luận tối đa & 5 \\
Số hành động giữ lại sau lọc (Top-K) & 12 \\
Trần ứng viên trước khi lọc & 80 \\
Kích thước tập kiểm thử (lấy mẫu phân tầng) & 200 \\
\bottomrule
\end{tabular}
\end{table}

Do tập kiểm thử lớn, chúng tôi áp dụng \textbf{lấy mẫu phân tầng} (\textit{stratified sampling}):
rút một tập con gồm 200 câu hỏi cân đối theo loại để đánh giá ở mỗi lượt, đúng theo quy trình của
bản gốc. Thiết lập phần cứng cụ thể (CPU/GPU, dung lượng bộ nhớ) và thời gian chạy trung bình mỗi
câu hỏi được ghi nhận trong Bảng~\ref{tab:setup} \todo{bổ sung thông tin phần cứng và thời gian chạy
khi đánh giá}.

\section{Các phương pháp so sánh và độ đo}
\label{sec:baselines}

\subsection{Các phương pháp so sánh}

Để làm rõ đóng góp của khung ARI, chúng tôi so sánh với một số phương pháp đại diện, được điều chỉnh
cho bối cảnh tiếng Việt và cùng dùng mô hình ngôn ngữ cục bộ để bảo đảm so sánh công bằng:

\begin{itemize}
  \item \textbf{LLM trực tiếp:} mô hình ngôn ngữ lớn trả lời thẳng câu hỏi mà không truy cập đồ thị
  tri thức. Phương pháp này đo năng lực suy luận thời gian dựa trên tri thức tham số sẵn có.
  \item \textbf{KG-RAG:} truy hồi các sự kiện liên quan (tối đa một số lượng nhất định) từ đồ thị
  theo thực thể và thời gian trong câu hỏi, đưa vào câu lệnh cho mô hình trả lời theo lối sinh tăng
  cường bằng truy hồi~\cite{baek2023kaping}.
  \item \textbf{CoT-KB:} kết hợp mô-đun tương tác tri thức vào khung chuỗi suy nghĩ
  (CoT)~\cite{wei2022cot}, cho mô hình tương tác với đồ thị dưới sự dẫn dắt của các ví dụ mẫu.
  \item \textbf{ReAct-KB:} biến thể của ReAct~\cite{yao2023react} tích hợp mô-đun không phụ thuộc
  tri thức (cùng tập thao tác nguyên tử như ARI); mô hình chọn một hành động trong các ứng viên ở
  mỗi bước, nhưng \emph{không} có hướng dẫn trừu tượng tùy biến theo câu hỏi.
  \item \textbf{ARI (đề xuất):} hệ thống đầy đủ của khóa luận với phương pháp luận trừu tượng.
\end{itemize}

\noindent Cả bốn phương pháp so sánh đều được cài đặt trong mã nguồn và chạy bằng cùng mô hình ngôn
ngữ cục bộ: ba baseline một lượt (\texttt{llm-only}, \texttt{kg-rag}, \texttt{cot-kb}) ở
\texttt{ari/baselines.py}, baseline nhiều bước \texttt{react-kb} dùng chung vòng lặp hành động với
ARI (\texttt{agent.run\_question(react=True)}). Tất cả gọi qua \texttt{evaluate.py -{}-method <tên>},
chấm điểm bằng cùng hàm so khớp đáp án như ARI để bảo đảm độ đo đồng nhất.

\subsection{Độ đo}

Độ đo chính là \textbf{độ chính xác} (\textit{accuracy}) theo Công thức~\eqref{eq:accuracy}: tỉ lệ
câu hỏi mà đáp án dự đoán nằm trong tập đáp án chuẩn. Kết quả được báo cáo ở ba mức: (i) tổng thể;
(ii) theo độ phức tạp -- nhóm \emph{đơn giản} (\texttt{simple\_entity}, \texttt{simple\_time}) và
nhóm \emph{phức tạp} (\texttt{before\_after}, \texttt{first\_last}, \texttt{time\_join}); và (iii)
theo kiểu đáp án (thực thể so với thời gian). Việc tách nhỏ này giúp làm rõ phương pháp mạnh ở dạng
suy luận nào.

\section{Kết quả tổng (CHNC1)}
\label{sec:overall-results}

Bảng~\ref{tab:overall} trình bày độ chính xác của ARI so với các phương pháp so sánh trên CronQ-VN.

\begin{table}[ht]
\centering
\caption{Độ chính xác trên CronQ-VN (đơn vị: \%). Các ô ``--'' sẽ được điền sau khi chạy đánh giá.}
\label{tab:overall}
\begin{tabular}{L{4.2cm} C{1.8cm} C{1.8cm} C{1.8cm} C{1.8cm} C{1.8cm}}
\toprule
\textbf{Phương pháp} & \textbf{Tổng} & \textbf{Đơn giản} & \textbf{Phức tạp} & \textbf{Thực thể} & \textbf{Thời gian} \\
\midrule
LLM trực tiếp & -- & -- & -- & -- & -- \\
KG-RAG & -- & -- & -- & -- & -- \\
CoT-KB & -- & -- & -- & -- & -- \\
ReAct-KB & -- & -- & -- & -- & -- \\
\textbf{ARI (đề xuất)} & -- & -- & -- & -- & -- \\
\bottomrule
\end{tabular}
\end{table}

\noindent\todo{điền số liệu Bảng~\ref{tab:overall} sau khi chạy \texttt{evaluate.py}; in đậm kết quả
tốt nhất mỗi cột.}

\medskip

Dựa trên kết quả của phương pháp gốc trên tiếng Anh~\cite{chen2024ari} và đặc thù của CronQ-VN,
chúng tôi \emph{kỳ vọng} các xu hướng sau, sẽ được kiểm chứng bằng số liệu thực tế:

\begin{itemize}
  \item \textbf{LLM trực tiếp đạt kết quả thấp}, đặc biệt với câu hỏi có đáp án là thời gian, do
  hạn chế về tri thức thời gian và suy luận nhiều bước (Mục~\ref{sec:limitations}).
  \item \textbf{Việc bổ sung tri thức (KG-RAG) cải thiện so với LLM trực tiếp} nhưng vẫn hạn chế do
  nhiễu và giới hạn ngữ cảnh.
  \item \textbf{ARI vượt trội ở nhóm câu hỏi phức tạp}, nơi hướng dẫn trừu tượng phát huy tác dụng
  trong việc định hướng chuỗi suy luận nhiều bước.
\end{itemize}

Bảng~\ref{tab:byqtype} trình bày chi tiết độ chính xác của ARI theo từng loại câu hỏi, giúp xác
định loại suy luận nào còn là thách thức.

\begin{table}[ht]
\centering
\caption{Độ chính xác của ARI theo từng loại câu hỏi trên CronQ-VN (đơn vị: \%).}
\label{tab:byqtype}
\begin{tabular}{L{4.5cm} C{3cm}}
\toprule
\textbf{Loại câu hỏi} & \textbf{Độ chính xác} \\
\midrule
\texttt{simple\_entity} & -- \\
\texttt{simple\_time}   & -- \\
\texttt{before\_after}  & -- \\
\texttt{first\_last}    & -- \\
\texttt{time\_join}     & -- \\
\bottomrule
\end{tabular}
\end{table}

\section{Nghiên cứu cắt bỏ (CHNC2)}
\label{sec:ablation}

Để đánh giá vai trò của từng thành phần trong khung ARI, chúng tôi tiến hành nghiên cứu cắt bỏ
(\textit{ablation study}), lần lượt loại bỏ một thành phần và đo mức suy giảm hiệu năng. Bốn biến
thể được xét, theo~\cite{chen2024ari}:

\begin{itemize}
  \item \textbf{Bỏ hướng dẫn trừu tượng} (\textit{w/o Abstract Guidance}): mô hình tự suy luận
  không có phương pháp luận. Biến thể này được hỗ trợ trực tiếp qua tùy chọn
  \texttt{-{}-no-methodology} trong \texttt{evaluate.py}.
  \item \textbf{Bỏ phân cụm lịch sử} (\textit{w/o History Cluster}): chắt lọc một phương pháp luận
  chung từ toàn bộ lịch sử, không phân theo loại câu hỏi. Tái tạo bằng \texttt{learn -{}-k 1
  -{}-resume} (dùng lại lịch sử đã chạy) rồi đánh giá với bộ nhớ một cụm này.
  \item \textbf{Bỏ lọc hành động} (\textit{w/o Action Filter}): cho mô hình chọn trong toàn bộ hành
  động sinh ra, không lọc Top-K. Hỗ trợ qua tùy chọn \texttt{-{}-no-action-filter} trong
  \texttt{evaluate.py}.
  \item \textbf{Bỏ ví dụ sai} (\textit{w/o Incorrect Examples}): chỉ dùng ví dụ đúng khi chắt lọc
  phương pháp luận. Tái tạo bằng \texttt{learn -{}-no-incorrect -{}-resume}.
\end{itemize}

\begin{table}[ht]
\centering
\caption{Kết quả nghiên cứu cắt bỏ trên CronQ-VN (độ chính xác tổng, \%).}
\label{tab:ablation}
\begin{tabular}{L{6.5cm} C{3cm}}
\toprule
\textbf{Cấu hình} & \textbf{Độ chính xác} \\
\midrule
ARI (đầy đủ) & -- \\
\quad bỏ hướng dẫn trừu tượng & -- \\
\quad bỏ phân cụm lịch sử & -- \\
\quad bỏ lọc hành động & -- \\
\quad bỏ ví dụ sai & -- \\
\bottomrule
\end{tabular}
\end{table}

\noindent\todo{điền số liệu Bảng~\ref{tab:ablation}.} Theo kết quả của bản gốc, dự kiến việc bỏ
hướng dẫn trừu tượng gây suy giảm mạnh nhất, khẳng định vai trò cốt lõi của phương pháp luận trừu
tượng; việc bỏ phân cụm cũng làm giảm hiệu năng, cho thấy một phương pháp luận chung là không đủ cho
các dạng câu hỏi khác nhau.

\section{Phân tích bộ lọc cross-encoder (CHNC3)}
\label{sec:ce-analysis}

Phần này đánh giá riêng đóng góp thứ ba: thay bộ lọc Top-K thuần cosine bằng bộ lọc hai tầng
cosine + cross-encoder (Mục~\ref{sec:crossencoder}). Chúng tôi so sánh hai cấu hình lọc trên cùng hệ
thống ARI, theo ba khía cạnh: độ chính xác cuối cùng, \emph{độ phủ hành động đúng} (tỉ lệ bước mà
hành động đúng được giữ lại sau lọc) và số bước suy luận trung bình.

\begin{table}[ht]
\centering
\caption{So sánh bộ lọc cosine và bộ lọc cross-encoder trên CronQ-VN.}
\label{tab:ce}
\begin{tabular}{L{5.0cm} C{2.6cm} C{2.6cm} C{2.6cm}}
\toprule
\textbf{Bộ lọc} & \textbf{Độ chính xác} & \textbf{Độ phủ hành động đúng} & \textbf{Số bước TB} \\
\midrule
Cosine Top-K (gốc) & -- & -- & -- \\
Cosine + Cross-encoder & -- & -- & -- \\
\bottomrule
\end{tabular}
\end{table}

\noindent\todo{điền số liệu Bảng~\ref{tab:ce} sau khi triển khai và huấn luyện cross-encoder.}

Giả thuyết kiểm chứng là: bằng cách xếp hạng hành động \emph{theo ngữ cảnh lịch sử} thay vì chỉ theo
câu hỏi gốc, bộ lọc cross-encoder giữ lại hành động đúng cho bước hiện tại với tỉ lệ cao hơn, qua đó
nâng độ chính xác cuối cùng, nhất là ở các câu hỏi phức tạp nhiều bước. Chúng tôi cũng dự kiến khảo
sát ảnh hưởng của kích thước tập ứng viên mà tầng cosine chọn ra trước khi đưa vào cross-encoder.

\section{Hiệu quả suy luận và phân tích lỗi}
\label{sec:error-analysis}

\subsection{Hiệu quả suy luận}

Ngoài độ chính xác, chúng tôi đánh giá \emph{hiệu quả suy luận} qua số bước trung bình mà hệ thống
cần để đi tới đáp án. Theo bản gốc~\cite{chen2024ari}, hướng dẫn trừu tượng không chỉ tăng độ chính
xác mà còn \emph{giảm} số bước suy luận, do mô hình được định hướng tốt hơn. Hình~\ref{fig:steps}
(khung tạm) so sánh số bước trung bình khi có và không có hướng dẫn trừu tượng; số liệu sẽ được điền
sau khi chạy.

\begin{figure}[ht]
  \centering
  \fbox{\parbox{0.8\textwidth}{\centering\vspace{6mm}
  \todo{biểu đồ cột: số bước suy luận trung bình, có vs không có hướng dẫn trừu tượng}
  \vspace{6mm}}}
  \caption{So sánh số bước suy luận trung bình khi có và không có hướng dẫn trừu tượng.}
  \label{fig:steps}
\end{figure}

Bên cạnh đó, chúng tôi khảo sát ảnh hưởng của \emph{số cụm} $k$ tới hiệu năng (Hình~\ref{fig:clusters},
khung tạm). Theo phân tích của bản gốc, hiệu năng tăng rồi giảm khi $k$ tăng: $k$ quá nhỏ khiến mỗi
cụm trộn nhiều dạng câu hỏi nên phương pháp luận kém sắc nét, còn $k$ quá lớn khiến mỗi cụm có quá ít
mẫu để chắt lọc phương pháp luận tốt; điểm tối ưu thường quanh $k = 10$.

\begin{figure}[ht]
  \centering
  \fbox{\parbox{0.8\textwidth}{\centering\vspace{6mm}
  \todo{biểu đồ đường: độ chính xác theo số cụm k (0--20)}
  \vspace{6mm}}}
  \caption{Ảnh hưởng của số cụm $k$ tới độ chính xác.}
  \label{fig:clusters}
\end{figure}

\subsection{Phân tích lỗi}
\label{sec:errors}

Phân tích định tính các trường hợp sai cho thấy ba nhóm lỗi điển hình, tương đồng với quan sát của
bản gốc~\cite{chen2024ari}:

\begin{enumerate}
  \item \textbf{Liên kết thực thể sai.} Hệ thống truy xuất nhầm thực thể từ đồ thị, thường do bước
  liên kết thực thể ánh xạ cụm từ trong câu hỏi về sai mã định danh. Lỗi này lan truyền và làm sai
  toàn bộ chuỗi suy luận phía sau.
  \item \textbf{Phương pháp luận trừu tượng kém chất lượng.} Với một số dạng câu hỏi phức tạp mà
  lịch sử suy luận hầu hết đều sai, mô hình không có đủ ví dụ đúng để chắt lọc phương pháp luận hiệu
  quả, dẫn tới hướng dẫn kém và gây ra chuỗi lỗi nối tiếp.
  \item \textbf{Đầu ra không tuân thủ định dạng.} Dù bị ràng buộc chỉ được chọn trong các hành động
  ứng viên hoặc đưa ra đáp án, đôi khi mô hình ngôn ngữ lớn không tuân thủ chặt chỉ dẫn, khiến truy
  vấn đồ thị thất bại và cản trở suy luận.
\end{enumerate}

Ba nhóm lỗi này gợi mở các hướng cải thiện: nâng cao chất lượng liên kết thực thể tiếng Việt, làm
giàu nguồn ví dụ đúng cho việc chắt lọc phương pháp luận, và tăng độ tuân thủ định dạng của mô hình
(ví dụ qua câu lệnh chặt chẽ hơn hoặc mô hình lớn hơn).

\section{Bàn luận và hạn chế}
\label{sec:discussion}

Kết quả và phân tích trên cho thấy khung ARI là một hướng tiếp cận phù hợp cho bài toán TKGQA tiếng
Việt, song vẫn còn một số hạn chế cần lưu ý.

Thứ nhất, chất lượng của phương pháp luận trừu tượng \emph{phụ thuộc nhiều vào năng lực của mô hình
ngôn ngữ lớn}. Khi vận hành cục bộ với mô hình quy mô vừa phải, chất lượng hướng dẫn có thể không
bằng các mô hình thương mại lớn, ảnh hưởng tới hiệu năng -- đây là đánh đổi giữa tính khả triển khai
cục bộ và độ mạnh của mô hình.

Thứ hai, ARI dựa trên \emph{suy luận nhiều bước}, nên thời gian suy luận dài hơn so với trả lời trực
tiếp; mỗi câu hỏi cần nhiều lần gọi mô hình, làm tăng chi phí tính toán.

Thứ ba, phạm vi của khóa luận tập trung vào suy luận \emph{thời gian} trên đồ thị tri thức; hiệu quả
của phương pháp trên các dạng suy luận khác chưa được khảo sát. Ngoài ra, bộ dữ liệu CronQ-VN hiện
giới hạn ở hạt thời gian cấp \emph{năm} và một tập quan hệ tuyển chọn, nên chưa bao quát hết độ phức
tạp của ngôn ngữ tự nhiên thực tế.

Những hạn chế này định hướng cho các đề xuất cải tiến và mở rộng trình bày ở Chương~\ref{ch:conclusion}.
